\chapter{Проектирование библиотеки}
\label{chap:library-design}

\section{Требования к библиотеке}

Проектирование библиотеки распределенной трассировки должно опираться на
практическую задачу, сформулированную в предыдущих главах: уменьшить объем ручной
интеграции средств наблюдаемости в HTTP-сервисах на Rust. Следовательно, при
проектировании необходимо выделить не только функциональные возможности, но и
нефункциональные требования, определяющие удобство внедрения и эксплуатационную
пригодность решения.

Функциональные требования описывают, какие возможности библиотека должна
предоставить прикладному сервису. Нефункциональные требования определяют
ограничения на способ использования этих возможностей: требование минимальной
конфигурации, совместимость с существующим стеком и предсказуемость поведения при
ошибках и завершении работы.

\begin{table}[H]
    \centering
    \caption{Функциональные требования к библиотеке}
    \label{tab:functional-requirements}
    \renewcommand{\arraystretch}{1.2}
    \begin{tabularx}{\textwidth}{>{\raggedright\arraybackslash}p{1.6cm} >{\raggedright\arraybackslash}X >{\raggedright\arraybackslash}p{2.1cm}}
        \toprule
        Обозначение & Содержание требования & Приоритет \\
        \midrule
        Ф1 & Библиотека должна инициализировать общую подсистему трассировки приложения на основе конфигурации сервиса & Обязательный \\
        \addlinespace
        Ф2 & Библиотека должна автоматически создавать серверный спан для каждого входящего HTTP-запроса & Обязательный \\
        \addlinespace
        Ф3 & Библиотека должна извлекать и внедрять контекст трассировки через W3C Trace Context & Обязательный \\
        \addlinespace
        Ф4 & Библиотека должна поддерживать трассировку исходящих HTTP-вызовов через специализированный клиентский интерфейс & Обязательный \\
        \addlinespace
        Ф5 & Библиотека должна поддерживать экспорт трейсов по OTLP/HTTP во внешнюю инфраструктуру наблюдаемости & Обязательный \\
        \addlinespace
        Ф6 & Библиотека должна добавлять идентификаторы трейса и спана в структурированные журналы событий & Обязательный \\
        \addlinespace
        Ф7 & Библиотека должна поддерживать безопасный захват разрешенных HTTP-заголовков в атрибуты спана & Желательный \\
        \addlinespace
        Ф8 & Библиотека должна поддерживать запись прикладных событий внутри активного трейса & Желательный \\
        \addlinespace
        Ф9 & Библиотека должна сопровождаться примерами интеграции с HTTP-сервисами на Axum & Желательный \\
        \bottomrule
    \end{tabularx}
\end{table}

\begin{table}[H]
    \centering
    \caption{Нефункциональные требования к библиотеке}
    \label{tab:nonfunctional-requirements}
    \renewcommand{\arraystretch}{1.2}
    \begin{tabularx}{\textwidth}{>{\raggedright\arraybackslash}p{1.6cm} >{\raggedright\arraybackslash}X >{\raggedright\arraybackslash}p{2.1cm}}
        \toprule
        Обозначение & Содержание требования & Приоритет \\
        \midrule
        Н1 & Подключение библиотеки должно требовать минимального количества пользовательского кода и параметров конфигурации & Обязательный \\
        \addlinespace
        Н2 & Библиотека должна быть совместима с экосистемой Tower и Axum~\cite{towerDocs,axumDocs} & Обязательный \\
        \addlinespace
        Н3 & Конфигурация должна проходить валидацию до запуска подсистемы трассировки & Обязательный \\
        \addlinespace
        Н4 & Обработка заголовков должна исключать запись чувствительных данных по умолчанию & Обязательный \\
        \addlinespace
        Н5 & Завершение работы экспортера должно быть управляемым и предсказуемым по времени & Желательный \\
        \addlinespace
        Н6 & Интерфейс библиотеки должен использовать понятные типы ошибок и устойчивые значения по умолчанию & Желательный \\
        \bottomrule
    \end{tabularx}
\end{table}

Перечисленные требования показывают, что библиотека должна решать не одну
локальную задачу, а целостную задачу внедрения трассировки в прикладной сервис.
Именно поэтому проектирование не может сводиться только к одному промежуточному
слою для входящих запросов или только к экспортеру. Требуется согласованная
архитектура, объединяющая конфигурацию, серверную обработку, клиентские вызовы,
журналирование и завершение работы.

Приоритеты требований связаны с минимально необходимой версией библиотеки.
Обязательные требования описывают путь, без которого распределенная трассировка
не складывается в рабочий сценарий: сервис должен инициализировать подсистему
трассировки, создавать серверные спаны, переносить контекст, инструментировать
исходящие вызовы, экспортировать трейсы и связывать журналы с активным контекстом.
Если убрать хотя бы один из этих элементов, пользователь снова будет вынужден
собирать часть инфраструктуры вручную.

Желательные требования имеют другую роль. Безопасный захват заголовков,
прикладные события и демонстрационные примеры не меняют базовую модель
трассировки, но делают библиотеку пригодной для реального использования. Захват
заголовков помогает добавить диагностический контекст, события позволяют
фиксировать доменные шаги внутри трейса, а примеры показывают повторяемый способ
подключения. Поэтому эти возможности проектируются как часть библиотеки, но не
подменяют ее основной контракт.

\section{Компонентная архитектура}

Верхнеуровневая архитектура библиотеки строится вокруг нескольких устойчивых
компонентов, каждый из которых отвечает за отдельную границу системы. С одной
стороны находятся прикладные HTTP-сервисы, использующие библиотеку. С другой
стороны находится внешняя инфраструктура наблюдаемости, принимающая и отображающая
трейсы. Между ними расположен библиотечный слой `my-infra-otel`, объединяющий
конфигурацию, обработку входящих и исходящих вызовов, распространение контекста,
экспорт и корреляцию журналов событий.

\begin{figure}[H]
    \centering
    \begin{tikzpicture}[
        node distance=1.8cm,
        block/.style={draw, rounded corners, minimum width=3.8cm, minimum height=0.95cm, align=center},
        arrow/.style={->, thick}
    ]
        \node[block] (apps) {Прикладные сервисы\\на Axum и Tower};
        \node[block, below=of apps] (sdk) {Библиотека `my-infra-otel`};
        \node[block, below left=1.6cm and 2.2cm of sdk] (layer) {Серверная обработка\\входящих запросов};
        \node[block, below=of sdk] (client) {Клиентская обработка\\исходящих запросов};
        \node[block, below right=1.6cm and 2.2cm of sdk] (logging) {Структурированные\\журналы событий};
        \node[block, below=3.0cm of client] (collector) {Сборщик телеметрии\\OTLP/HTTP};
        \node[block, below=of collector] (jaeger) {Jaeger};

        \draw[arrow] (apps) -- (sdk);
        \draw[arrow] (sdk) -- (layer);
        \draw[arrow] (sdk) -- (client);
        \draw[arrow] (sdk) -- (logging);
        \draw[arrow] (sdk) -- (collector);
        \draw[arrow] (collector) -- (jaeger);
    \end{tikzpicture}
    \caption{Компонентная архитектура библиотеки распределенной трассировки}
    \label{fig:component-architecture}
\end{figure}

На рисунке~\ref{fig:component-architecture} показано, что библиотека играет роль
связующего слоя между прикладным кодом и инфраструктурой наблюдаемости. Она не
хранит трейсы самостоятельно и не подменяет прикладную бизнес-логику, а
обеспечивает единый путь подключения трассировки на границе входящих и исходящих
вызовов. Такое разделение обязанностей соответствует технической схеме,
зафиксированной в `ARCHITECTURE.md`, и позволяет проектировать библиотеку как
самостоятельный слой интеграции, а не как часть конкретного сервиса.

С точки зрения потоков зависимостей прикладные сервисы зависят от публичного
интерфейса библиотеки, библиотека зависит от `tracing`, OpenTelemetry, Tower,
Axum и `reqwest`, а экспорт трейсов направляется во внешнюю систему через протокол
OTLP/HTTP~\cite{otlpSpec,otlpExporterSpec}. Благодаря этому прикладной код не
содержит прямой зависимости от Jaeger и не должен самостоятельно собирать
низкоуровневую цепочку экспорта.

Границы ответственности в этой архитектуре разделены явно. Прикладной сервис
отвечает за бизнес-логику, маршруты, обработчики и выбор зависимых вызовов.
Библиотека отвечает за единообразное создание спанов, перенос контекста,
структурированные журналы и настройку экспорта. OpenTelemetry Collector принимает
телеметрию по OTLP/HTTP, выполняет промежуточную обработку и передает данные
дальше. Jaeger отвечает за хранение и визуальное представление трейсов. Такое
разделение снижает связность: сервисы не знают деталей Jaeger, а библиотека не
берет на себя роль хранилища или пользовательского интерфейса наблюдаемости.

Важно также, что внешний контур наблюдаемости проектируется как заменяемый.
Использование OTLP/HTTP означает, что библиотека передает данные в стандартный
приемник, а не в конкретную реализацию хранилища. В демонстрационном контуре
таким приемником является Collector, а системой просмотра --- Jaeger, однако
сам публичный интерфейс библиотеки не требует от прикладного кода прямого знания
об этих компонентах.

\section{Модули библиотеки}

Проектирование библиотечного слоя реализовано через набор модулей с четко
разделенной ответственностью. Такое разбиение позволяет сохранить понятный
публичный интерфейс и изолировать инфраструктурные детали в отдельных частях
системы.

\begin{table}[H]
    \centering
    \caption{Ответственность основных модулей библиотеки}
    \label{tab:module-responsibilities}
    \renewcommand{\arraystretch}{1.2}
    \begin{tabularx}{\textwidth}{>{\raggedright\arraybackslash}p{3.0cm} >{\raggedright\arraybackslash}X >{\raggedright\arraybackslash}X}
        \toprule
        Модуль & Основная ответственность & Взаимодействие \\
        \midrule
        `config` & Описание конфигурации сервиса, значения по умолчанию и валидация параметров & Используется в `init`, `logging`, `internal::otel`, `guard` \\
        \addlinespace
        `layer` & Создание серверного спана для входящих HTTP-запросов и фиксация их атрибутов & Использует `propagation`, `header_capture`, `tracing` \\
        \addlinespace
        `propagation` & Извлечение и внедрение контекста трассировки через HTTP-заголовки & Используется в `layer` и `request_builder` \\
        \addlinespace
        `internal::otel` и `init` & Настройка поставщика трассировщиков, распространителя контекста и экспортера OTLP & Используют `TracingConfig`, OpenTelemetry и `tracing-subscriber` \\
        \addlinespace
        `logging` & Формирование структурированных журналов и добавление идентификаторов трейса и спана & Использует данные конфигурации и активный контекст трассировки \\
        \addlinespace
        `client` и `request_builder` & Построение и отправка трассируемых исходящих HTTP-запросов & Используют `reqwest`, `propagation`, `tracing` \\
        \addlinespace
        `events` & Запись прикладных событий в рамках активного трейса & Использует структурированные поля событий \\
        \addlinespace
        `error` & Единые типы ошибок конфигурации и выполнения & Используется всеми публичными интерфейсами \\
        \addlinespace
        `lib` & Переэкспорт публичных сущностей и точка входа в библиотеку & Связывает внутренние модули с прикладным кодом \\
        \bottomrule
    \end{tabularx}
\end{table}

\subsection{Модуль конфигурации}

Модуль `config` определяет структуру `TracingConfig` и связанный с ней построитель
`TracingConfigBuilder`. Проектное решение состоит в том, чтобы вынести все
существенные параметры в одну конфигурационную модель: имя сервиса, версию,
окружение, адрес конечной точки OTLP, фильтр журналирования, пользовательские
атрибуты ресурса, режим выборки, тайм-аут экспорта, тайм-аут завершения и формат
журналов. Благодаря этому инициализация трассировки получает единый входной
контракт.

Построитель задает также значения по умолчанию. В частности, окружение получает
значение `local`, адрес экспорта по умолчанию указывает на локальный приемник
OTLP/HTTP, а режим выборки по умолчанию всегда включен. Важным проектным решением
является встроенная валидация конфигурации до запуска подсистемы трассировки.
Таким образом, модуль `config` отвечает не только за хранение параметров, но и за
контроль допустимости начального состояния системы.

\subsection{Серверная обработка входящих запросов}

Модуль `layer` реализует промежуточный слой `MyOtelTracingLayer` и соответствующий
построитель. Этот модуль отвечает за границу входа HTTP-запроса в сервис. При
проектировании было важно использовать стандартные абстракции Tower `Layer` и
`Service`, поскольку именно они обеспечивают совместимость с Axum и позволяют
автоматически применять библиотеку ко всем маршрутам сервиса~\cite{towerDocs,axumDocs}.

При вызове внутреннего сервиса слой должен определить метод и путь запроса,
получить шаблон маршрута, создать серверный спан `http.request`, извлечь
родительский контекст из заголовков и затем выполнить обработчик внутри этого
спана. После завершения обработки он обязан зафиксировать код ответа,
длительность и признаки ошибки. Такое поведение делает модуль `layer` центральным
элементом проектной схемы серверной трассировки.

\subsection{Передача контекста}

Модуль `propagation` реализует извлечение и внедрение контекста трассировки через
HTTP-заголовки. Проектное решение здесь состоит в том, чтобы не смешивать логику
обработки запроса с деталями OpenTelemetry. Для этого используются отдельные
адаптеры над `HeaderMap`, которые преобразуют набор заголовков к интерфейсам
извлечения и внедрения. В результате модуль имеет узкую и четкую ответственность:
переносить контекст трассировки между процессами в соответствии с W3C Trace
Context~\cite{w3cTraceContext,opentelemetryPropagationDocs}.

\subsection{Экспорт и инициализация}

Модули `init`, `internal::otel` и `guard` образуют группу, отвечающую за запуск и
завершение подсистемы трассировки. `init_global_tracing` выполняет однократную
инициализацию общей подсистемы и защищает приложение от повторного запуска через
атомарный флаг. Внутренний модуль `internal::otel` создает поставщик
трассировщиков, настраивает распространитель контекста W3C, подготавливает
ресурсные атрибуты сервиса и собирает конвейер экспорта через OTLP/HTTP. Модуль
`guard` обеспечивает управляемое завершение работы экспортера и предоставляет
интерфейс `shutdown`, ограниченный тайм-аутом.

Такое разбиение важно, поскольку оно изолирует инфраструктурные детали от
прикладного кода. Прикладной сервис вызывает один публичный метод и получает один
объект управления жизненным циклом, не взаимодействуя напрямую с низкоуровневыми
компонентами OpenTelemetry.

\subsection{Журналы событий}

Модуль `logging` реализует форматтер `JsonTraceFormatter`, который обогащает
структурированные журналы сведениями о сервисе и активном трейсе. В проектной
схеме это необходимо для корреляции логов и трейсов: если событие возникает внутри
активного спана, форматтер добавляет идентификатор трейса и идентификатор спана.
Таким образом, журналы остаются полезными для прикладной диагностики, но
одновременно связываются с общей системой наблюдаемости.

\subsection{Клиентский интерфейс и прикладные события}

Модули `client` и `request_builder` отвечают за исходящие HTTP-вызовы.
Проектирование отдельного клиентского интерфейса `TracedHttpClient` позволяет не
вмешиваться в низкоуровневую реализацию `reqwest::Client`, а обернуть ее
специализированным интерфейсом, который автоматически создает клиентский спан,
внедряет контекст и фиксирует результат выполнения вызова.

Модуль `events` дополняет проектный слой поддержкой прикладных событий. Функция
`record_event` принимает имя события и набор структурированных полей. Это решение
позволяет добавлять доменный контекст в трейс без ручной сборки низкоуровневых
записей и без смешивания бизнес-логики с деталями OpenTelemetry.

\subsection{Публичный интерфейс и ошибки}

Корневой модуль `lib` проектируется как единая точка входа в библиотеку. Он
переэкспортирует конфигурацию, промежуточный слой, клиентские типы, события,
политику захвата заголовков и типы ошибок. За счет этого прикладной сервис
использует небольшое число устойчивых публичных сущностей и не зависит от
внутренней структуры модулей.

Модуль `error` задает две основные группы ошибок: ошибки конфигурации
`ConfigError` и ошибки выполнения `MyOtelError`. Такое разделение повышает
предсказуемость публичного интерфейса: потребитель библиотеки может отличать
ошибки пользовательской настройки от ошибок инициализации экспортера,
журналирования или сетевого клиента.

\subsection{Жизненный цикл и обработка ошибок}

Жизненный цикл библиотеки начинается с построения `TracingConfig` и вызова
`init_global_tracing`. На этом этапе выполняется проверка конфигурации,
создается поставщик трассировщиков, настраивается подписчик `tracing` и
подключается формат журналов. Проектное решение состоит в том, что прикладной
код получает не набор разрозненных объектов, а один объект `TracingGuard`,
который представляет запущенную подсистему трассировки.

Отдельный объект жизненного цикла нужен для управляемого завершения. Экспортер
может накапливать данные перед отправкой, поэтому при остановке процесса важно
дать ему возможность выполнить сброс. Метод `TracingGuard::shutdown()` задает
явную точку завершения, а тайм-аут в конфигурации ограничивает время ожидания.
Таким образом, завершение становится частью публичного контракта, а не
неявной побочной операцией.

Ошибки также проектируются как часть контракта. `ConfigError` описывает ошибки,
которые можно обнаружить до запуска: пустое имя сервиса, некорректную конечную
точку OTLP, недопустимый ключ атрибута, запрещенный заголовок или нулевой
тайм-аут. `MyOtelError` объединяет эти ошибки с ошибками инициализации,
клиентского HTTP-вызова, завершения работы и записи событий. Такое разделение
позволяет прикладному сервису реагировать на сбои предсказуемо: ошибка
конфигурации требует исправления настроек, а ошибка выполнения относится уже к
работе инфраструктуры или зависимого вызова.

\section{Публичный интерфейс библиотеки}

Публичный интерфейс должен быть достаточно компактным, чтобы библиотека
поддерживала низкопороговое подключение, и одновременно достаточно выразительным,
чтобы покрывать основные сценарии настройки и использования. В текущем проекте
этот интерфейс строится вокруг трех групп сущностей: конфигурации, промежуточного
слоя для входящих запросов и клиентского интерфейса для исходящих запросов.

\begin{table}[H]
    \centering
    \caption{Основные элементы публичного интерфейса}
    \label{tab:public-api}
    \renewcommand{\arraystretch}{1.2}
    \begin{tabularx}{\textwidth}{>{\raggedright\arraybackslash}p{4.0cm} >{\raggedright\arraybackslash}X >{\raggedright\arraybackslash}p{3.1cm}}
        \toprule
        Элемент интерфейса & Назначение & Результат использования \\
        \midrule
        `TracingConfig::builder(...)` & Начало построения конфигурации трассировки сервиса & `TracingConfigBuilder` \\
        \addlinespace
        `init_global_tracing(config)` & Инициализация общей подсистемы трассировки & `TracingGuard` \\
        \addlinespace
        `MyOtelTracingLayer::builder()` & Настройка промежуточного слоя для входящих HTTP-запросов & `MyOtelTracingLayerBuilder` \\
        \addlinespace
        `HeaderCapturePolicy::builder()` & Настройка правил захвата разрешенных заголовков & `HeaderCapturePolicyBuilder` \\
        \addlinespace
        `TracedHttpClient::new(client)` & Оборачивание клиентского HTTP-интерфейса трассируемой логикой & `TracedHttpClient` \\
        \addlinespace
        `TracedHttpClient::get`, `post`, `request` & Построение трассируемого исходящего запроса & `TracedRequestBuilder` \\
        \addlinespace
        `record_event(...)` & Запись прикладного события внутри активного трейса & Информационная запись трассировки \\
        \addlinespace
        `TracingGuard::shutdown()` & Управляемое завершение и сброс экспортера & `Result<()>` \\
        \bottomrule
    \end{tabularx}
\end{table}

Основной путь подключения библиотеки к сервису выглядит следующим образом:
\begin{enumerate}
    \item создать конфигурацию сервиса через построитель;
    \item инициализировать общую подсистему трассировки;
    \item построить промежуточный слой для входящих запросов;
    \item добавить этот слой к маршрутизатору HTTP-сервиса;
    \item использовать специализированный клиентский интерфейс для зависимых вызовов;
    \item при необходимости записывать прикладные события через `record_event`.
\end{enumerate}

Ниже приведен сокращенный пример такого подключения:

\begin{verbatim}
let config = TracingConfig::builder("checkout-api")
    .environment("local")
    .build()?;

let _guard = init_global_tracing(config)?;

let layer = MyOtelTracingLayer::builder()
    .headers(HeaderCapturePolicy::builder()
        .standard_request_ids()
        .build()?)
    .build()?;
\end{verbatim}

Пример показывает, что публичный интерфейс проектируется вокруг компактной
последовательности действий. Пользователь библиотеки не должен вручную создавать
поставщик трассировщиков, настраивать распространитель контекста или собирать
серверные и клиентские спаны на каждой границе вызова.

\begin{table}[H]
    \centering
    \caption{Типовые сценарии использования публичного интерфейса}
    \label{tab:api-scenarios}
    \renewcommand{\arraystretch}{1.2}
    \begin{tabularx}{\textwidth}{>{\raggedright\arraybackslash}p{3.5cm} >{\raggedright\arraybackslash}X >{\raggedright\arraybackslash}X}
        \toprule
        Сценарий & Используемые элементы API & Результат для прикладного сервиса \\
        \midrule
        Минимальное подключение трассировки & `TracingConfig::builder`, `init_global_tracing`, `MyOtelTracingLayer::new` & Сервис получает серверные спаны входящих HTTP-запросов и базовый экспорт трейсов \\
        \addlinespace
        Подключение с политикой заголовков & `HeaderCapturePolicy::builder`, `MyOtelTracingLayer::builder` & В спаны попадают только разрешенные диагностические заголовки с заданным режимом записи \\
        \addlinespace
        Исходящий HTTP-вызов & `TracedHttpClient::new`, `get`, `post`, `request` & Исходящий запрос получает клиентский спан и переносит текущий контекст в заголовки \\
        \addlinespace
        Запись прикладного события & `record_event`, `EventField`, `EventValue` & В активном трейсе появляется доменное событие со структурированными полями \\
        \addlinespace
        Управляемое завершение & `TracingGuard::shutdown` & Сервис явно завершает подсистему трассировки и инициирует сброс накопленных данных \\
        \bottomrule
    \end{tabularx}
\end{table}

Таблица~\ref{tab:api-scenarios} показывает, что публичный интерфейс покрывает не
одну точку входа, а весь жизненный путь использования библиотеки. Сначала
пользователь настраивает и запускает подсистему трассировки, затем подключает
серверную обработку, после этого использует клиентский интерфейс для зависимых
вызовов и при необходимости дополняет трейс прикладными событиями. Завершение
работы также остается явным действием, что делает поведение библиотеки
предсказуемым при остановке сервиса.

\section{Потоки обработки запросов}

Последним этапом проектирования является фиксация типовых потоков обработки
входящих и исходящих HTTP-запросов. Эти потоки определяют, как взаимодействуют
между собой модули библиотеки и какие точки расширения получает прикладной сервис.

\subsection{Поток входящего запроса}

\begin{figure}[H]
    \centering
    \begin{tikzpicture}[
        node distance=1.5cm,
        block/.style={draw, rounded corners, minimum width=4.3cm, minimum height=0.9cm, align=center},
        arrow/.style={->, thick}
    ]
        \node[block] (request) {Входящий HTTP-запрос};
        \node[block, below=of request] (extract) {Извлечение контекста\\из заголовков};
        \node[block, below=of extract] (span) {Создание серверного спана\\`http.request`};
        \node[block, below=of span] (handler) {Выполнение обработчика\\внутри спана};
        \node[block, below=of handler] (finish) {Запись кода ответа,\\длительности и ошибок};

        \draw[arrow] (request) -- (extract);
        \draw[arrow] (extract) -- (span);
        \draw[arrow] (span) -- (handler);
        \draw[arrow] (handler) -- (finish);
    \end{tikzpicture}
    \caption{Поток обработки входящего HTTP-запроса}
    \label{fig:incoming-flow}
\end{figure}

Поток входящего запроса начинается на серверной границе. Промежуточный слой
получает HTTP-запрос, извлекает контекст трассировки из заголовков, создает
серверный спан и затем выполняет внутренний обработчик в рамках этого спана.
После завершения обработки слой записывает код ответа, длительность и признаки
ошибки. Если обработчик внутри себя инициирует дополнительные операции, они
становятся дочерними относительно серверного спана.

\subsection{Поток исходящего запроса}

\begin{figure}[H]
    \centering
    \begin{tikzpicture}[
        node distance=1.5cm,
        block/.style={draw, rounded corners, minimum width=4.5cm, minimum height=0.9cm, align=center},
        arrow/.style={->, thick}
    ]
        \node[block] (build) {Построение исходящего\\HTTP-запроса};
        \node[block, below=of build] (clientspan) {Создание клиентского спана\\`http.client.request`};
        \node[block, below=of clientspan] (inject) {Внедрение контекста\\в заголовки};
        \node[block, below=of inject] (send) {Отправка запроса\\в зависимый сервис};
        \node[block, below=of send] (result) {Запись результата,\\длительности и ошибок};

        \draw[arrow] (build) -- (clientspan);
        \draw[arrow] (clientspan) -- (inject);
        \draw[arrow] (inject) -- (send);
        \draw[arrow] (send) -- (result);
    \end{tikzpicture}
    \caption{Поток обработки исходящего HTTP-запроса}
    \label{fig:outgoing-flow}
\end{figure}

Поток исходящего запроса начинается в момент, когда прикладной код использует
специализированный клиентский интерфейс библиотеки. Построитель исходящего запроса
создает клиентский спан, внедряет в заголовки текущий контекст трассировки и
только после этого выполняет вызов зависимого сервиса. По завершении вызова в
спан записываются код ответа, длительность и признаки ошибки. Благодаря этому
исходящий запрос логически связывается как с текущим сервисом, так и с
последующим продолжением трейса в зависимой системе.

Таким образом, проектирование библиотеки завершает формирование целостного
архитектурного контракта. Требования задают ожидаемое поведение, компоненты и
модули определяют распределение ответственности, публичный интерфейс фиксирует
способ подключения, а потоки запросов показывают, как эти элементы работают в
типичном распределенном HTTP-сценарии. Это создает достаточную основу для
следующей главы, посвященной реализации и тестированию библиотеки.
